\chapter{Hypotheses}

Based on preliminary results and the literature, the central hypothesis of this work is that the computational model will reveal
a functional dichotomy between Granule Cell populations. It is expected that mGCs will behave as pattern separators, increasing
the dissimilarity between the neural representations of distinct stimuli. In contrast, iGCs, by virtue of their greater
intrinsic excitability, are predicted to function as pattern integrators, with their non-selective activation decreasing the
distinction between output patterns and, consequently, degrading the overall pattern separation capacity of the DG.

Additionally, it is hypothesized that this indiscriminate activation of iGCs will negatively impact CA3 memory functions,
impairing the fidelity of pattern completion by co-activating neural assemblies unrelated to the input stimulus. Although iGC
activity may suppress that of mGCs through shared inhibitory circuits, it is expected that this suppression will marginally
increase mGC sparsity, slightly improving their separation capacity. Finally, the model will test the hypothesis that iGCs act
as temporal integrators, with their maturation specializing them to respond to information encoded during their immature phase,
establishing a mechanism for the separation of memories in time.
